\documentclass[11pt,a4paper]{article}
\usepackage{fontspec}
\setmainfont{Times New Roman}
\setsansfont{Arial}
\usepackage[russian]{babel}
\usepackage{amsmath,amssymb}
\usepackage[margin=19mm,headheight=15pt,headsep=7mm,footskip=10mm]{geometry}
\usepackage{xcolor,fancyhdr,titlesec,enumitem,hyperref}
\definecolor{muted}{HTML}{000000}
\hypersetup{unicode=true,colorlinks=true,linkcolor=black,pdftitle={Госэкзамен: главные теоремы программы и их доказательства}}
\pagestyle{fancy}\fancyhf{}
\fancyhead[L]{\small\sffamily\color{muted} Госэкзамен · Главные теоремы программы}
\fancyfoot[L]{\footnotesize\color{muted} По темам программы · На основе конспекта 2021}
\fancyfoot[R]{\color{muted}\thepage}
\renewcommand{\headrulewidth}{0pt}
\setlength{\parindent}{0pt}\setlength{\parskip}{6pt}
\setlength{\abovedisplayskip}{7pt}\setlength{\belowdisplayskip}{7pt}
\setlist[enumerate]{leftmargin=19pt,itemsep=5pt,topsep=3pt,parsep=2pt}
\titleformat{\section}{\large\bfseries}{}{0pt}{}
\titlespacing*{\section}{0pt}{0pt}{10pt}
\emergencystretch=3em
\newcommand{\R}{\mathbb R}\newcommand{\C}{\mathbb C}
\newcommand{\E}{\mathbb E}\newcommand{\Var}{\operatorname{Var}}
\newcommand{\topic}[3]{\clearpage\section*{#1. #2}\phantomsection\addcontentsline{toc}{section}{#1. #2}{\small\color{muted}Конспект: #3.}\par}
\newcommand{\claim}{\textbf{Утверждение.} }
\newcommand{\idea}[1]{\par\textbf{Идея.} #1\par}
\newcommand{\uses}[1]{{\small\color{muted}\textbf{Используем:} #1.}\par}
\newcommand{\finish}{\hfill$\square$\par}
\newcommand{\remember}[1]{\par\textbf{Что запомнить.} #1}

\newcommand{\program}[1]{\par\begingroup\small\textbf{Пункт программы:} #1\par\endgroup\vspace{5pt}}



\begin{document}
\pagecolor{white}\color{black}
\section*{Главные теоремы программы: доказательства}
\textbf{38 разделов, сопоставленных с исходной программой госэкзамена.}

В начале каждого раздела приведена формулировка соответствующего пункта программы. Заголовок раздела показывает, какая именно главная теорема или основное утверждение выбраны для доказательства. В некоторых пунктах несколько тесно связанных результатов доказаны вместе: например, Чебышёв и закон больших чисел.

Приоритет отдан содержанию названия пункта. В частности, в пунктах 14--16 доказываются сами интегральные формулы, в 24--25 --- критерий диагонализируемости и спектральная теорема, в 27 --- метод для правой части-квазимногочлена, в 37 --- сама теорема Лорана.

\textbf{Как читать.} Сначала условия и утверждение, затем идея, после неё нумерованные шаги. Строка «Используем» перечисляет опорные результаты, принимаемые известными. В конце выделено, что важно не перепутать при устном ответе.

\textbf{О полноте доказательств.} Доказательства главных утверждений даны с указанными опорными теоремами, а не с выводом всей математики с нуля. В ЦПТ используется теорема Леви; в Лоране и вычетах --- интегральная теорема Коши. Для Грина и Гаусса используется геометрическая локализация, для Стокса --- разбиение поверхности на параметрические куски. Эти факты названы явно.

Разбиение единицы в разделах 14--15 означает выбор конечного набора гладких весов $\chi_j$, каждый из которых равен нулю вне своего малого участка, и $\sum_j\chi_j=1$ около рассматриваемой области. Тогда поле равно сумме локальных полей $\chi_jF$. Носитель функции --- замыкание множества, где она ненулевая.

\textbf{Нумерация.} Пункты 1--38 взяты из программы 2025--2026, ссылки на страницы и параграфы относятся к конспекту \textit{2021\_Oral (1) (2).pdf}. Переносы строк программы объединены; нумерация сохранена. Пункт 7 программы --- именно теорема о равномерной непрерывности на компакте, то есть теорема Кантора.

Это подборка главных доказательств по каждому пункту, а не полный набор доказательств всех результатов каждого билета.
\clearpage
{\small\tableofcontents}


\topic{1}{Больцано--Вейерштрасс}{§ 1.1, с. 7--9}
\program{Теорема Больцано-Вейерштрасса и критерий Коши сходимости числовой последовательности.}
\claim Из любой ограниченной вещественной последовательности $(x_n)$ можно выделить сходящуюся подпоследовательность.
\idea{Оставлять половину отрезка, содержащую бесконечно много членов. Длины отрезков стремятся к нулю.}
\uses{теорему о вложенных замкнутых отрезках}
\begin{enumerate}
\item Все $x_n$ лежат в некотором $I_0=[a,b]$. Если $a=b$, результат очевиден. Далее $a<b$.
\item Разделим $I_0$ пополам. Хотя бы одна половина содержит $x_n$ для бесконечно многих индексов: иначе их объединение содержало бы лишь конечное число индексов. Эту половину назовём $I_1$.
\item Повторяя выбор, получим замкнутые отрезки
\[
 I_0\supset I_1\supset I_2\supset\cdots,\qquad |I_k|=(b-a)2^{-k}.
\]
Существует единственная точка $c$, принадлежащая всем $I_k$.
\item Выбираем $n_k>n_{k-1}$ так, чтобы $x_{n_k}\in I_k$. Это возможно: подходящих индексов бесконечно много.
\item Обе точки $c$ и $x_{n_k}$ лежат в $I_k$, поэтому
\[
 |x_{n_k}-c|\le(b-a)2^{-k}\longrightarrow0.
\]
Следовательно, $x_{n_k}\to c$.
\end{enumerate}\finish
\remember{Индексы выбранных членов обязательно должны строго возрастать.}

\topic{2}{Достижение максимума на отрезке}{§ 1.2, с. 9--10}
\program{Ограниченность функции, непрерывной на отрезке, достижение точных верхней и нижней граней.}
\claim Непрерывная на $[a,b]$ вещественная функция ограничена и достигает максимума и минимума.
\idea{Выделить сходящуюся подпоследовательность точек и применить непрерывность.}
\uses{Больцано--Вейерштрасса; непрерывность; свойство супремума}
\begin{enumerate}
\item Если $f$ не ограничена, для каждого $n$ найдётся $x_n\in[a,b]$ с $|f(x_n)|>n$. Выделим $x_{n_k}\to c\in[a,b]$; принадлежность предела следует из замкнутости отрезка.
\item По непрерывности $f(x_{n_k})\to f(c)$, поэтому эта последовательность ограничена. Но $|f(x_{n_k})|>n_k\to\infty$. Противоречие: ограниченность доказана.
\item Существует конечное $M=\sup f([a,b])$. По свойству супремума выберем $y_n\in[a,b]$ так, что
\[
 M-1/n<f(y_n)\le M.
\]
Отсюда $f(y_n)\to M$.
\item Выделим $y_{n_k}\to d\in[a,b]$. Непрерывность даёт
\[
 f(d)=\lim_{k\to\infty}f(y_{n_k})=M.
\]
Максимум достигнут. Применяя доказанное к $-f$, получаем достижение минимума $f$.
\end{enumerate}\finish
\remember{Супремум --- число; для максимума нужна точка, в которой это число является значением функции.}

\topic{3}{Теорема о промежуточном значении}{§ 1.3, с. 11}
\program{Теорема о промежуточных значениях непрерывной функции.}
\claim Если $f\in C[a,b]$ и $f(a)<c<f(b)$, то существует $\xi\in(a,b)$ с $f(\xi)=c$. При обратном порядке значений вывод тот же.
\idea{Для $g=f-c$ всё время сохранять разные знаки на концах отрезка.}
\uses{теорему о вложенных отрезках; непрерывность}
\begin{enumerate}
\item Для $g(x)=f(x)-c$ имеем $g(a)<0<g(b)$.
\item Делим отрезок пополам. Если в середине $g=0$, всё доказано. Иначе выбираем половину, на левом конце которой $g<0$, а на правом $g>0$.
\item Повторяем, пока нуль не найден. Если процедура бесконечна, получаем $[a_n,b_n]$ с
\[
 g(a_n)<0<g(b_n),\qquad b_n-a_n=(b-a)2^{-n}.
\]
По теореме о вложенных отрезках $a_n\to\xi$ и $b_n\to\xi$ для одной точки $\xi$.
\item По непрерывности обе последовательности значений стремятся к $g(\xi)$. Предел отрицательных чисел не больше нуля, а положительных --- не меньше нуля. Значит, $g(\xi)\le0$ и $g(\xi)\ge0$, откуда $g(\xi)=0$.
\item Итак, $f(\xi)=c$. Поскольку $g(a)$ и $g(b)$ ненулевые, точка $\xi$ внутренняя. Обратный порядок знаков сводится к этому заменой $g$ на $-g$.
\end{enumerate}\finish
\remember{Без непрерывности функция может «перепрыгнуть» нужное значение.}

\topic{4}{Теорема Лагранжа}{§ 1.4.2, с. 12--13}
\program{Теоремы о среднем Ролля, Лагранжа и Коши для дифференцируемых функций.}
\claim Если $f\in C[a,b]$, $a<b$, и $f$ дифференцируема на $(a,b)$, то для некоторой $\xi\in(a,b)$
\[
 f(b)-f(a)=f'(\xi)(b-a).
\]
\idea{Вычесть прямую, соединяющую концы графика. У разности значения на концах равны.}
\uses{теоремы Вейерштрасса и Ферма}
\begin{enumerate}
\item Положим $k=(f(b)-f(a))/(b-a)$ и
\[
 h(x)=f(x)-f(a)-k(x-a).
\]
Тогда $h\in C[a,b]$, дифференцируема внутри и $h(a)=h(b)=0$.
\item Если $h\equiv0$, то $h'=0$ всюду внутри; можно выбрать любую внутреннюю точку.
\item Иначе $h$ принимает положительное или отрицательное значение. В первом случае её максимум положителен, во втором минимум отрицателен. Этот экстремум достигается по Вейерштрассу и не может лежать на концах, где $h=0$.
\item Значит, во внутренней точке $\xi$ по Ферма $h'(\xi)=0$. Так как $h'=f'-k$, имеем
\[
 f'(\xi)=\frac{f(b)-f(a)}{b-a}.
\]
\end{enumerate}\finish
\remember{Шаги 2--3 одновременно доказывают теорему Ролля для вспомогательной функции.}

\topic{5}{Тейлор с остатком Лагранжа}{§ 1.5.3, с. 15--17}
\program{Формула Тейлора с остаточным членом в форме Пеано или Лагранжа.}
\claim Пусть $n\ge0$, $f\in C^n$ на отрезке между $a$ и $x\ne a$, а $f^{(n+1)}$ существует внутри. Для некоторой точки $\xi$ строго между $a$ и $x$
\[
 f(x)=\sum_{j=0}^n\frac{f^{(j)}(a)}{j!}(x-a)^j
 +\frac{f^{(n+1)}(\xi)}{(n+1)!}(x-a)^{n+1}.
\]
\idea{Подобрать последний коэффициент так, чтобы многократно применить Ролля.}
\uses{теорему Ролля}
\begin{enumerate}
\item Обозначим многочлен в сумме через $T_n$. Зафиксируем $x$ и положим
\[
 K=\frac{f(x)-T_n(x)}{(x-a)^{n+1}},\quad
 H(t)=f(t)-T_n(t)-K(t-a)^{n+1}.
\]
Имеем $H(x)=0$ и $H(a)=H'(a)=\cdots=H^{(n)}(a)=0$.
\item Из $H(a)=H(x)=0$ по Роллю получаем внутренний нуль $H'$. Если $n\ge1$, у $H'$ есть также нуль в $a$, поэтому между ними найдётся нуль $H''$.
\item Продолжая, после $n+1$ применений получаем $H^{(n+1)}(\xi)=0$ между $a$ и $x$. Гладкость в условии обеспечивает все применения Ролля. При $n=0$ достаточно первого.
\item Так как $T_n^{(n+1)}=0$, имеем
\[
 0=f^{(n+1)}(\xi)-K(n+1)!.
\]
Полученное $K=f^{(n+1)}(\xi)/(n+1)!$ подставляем в его определение.
\end{enumerate}\finish
\remember{$x$ зафиксирован; $K$ зависит от этого $x$, но не от переменной $t$.}

\topic{6}{Необходимое и достаточное условия экстремума}{§ 1.6.2, с. 17--18}
\program{Исследование функций одной переменной при помощи первой и второй производных на монотонность, локальные экстремумы, выпуклость. Необходимые условия, достаточные условия.}
\claim Пусть $a$ --- внутренняя точка интервала определения $f$.
\textbf{Необходимое условие:} если в $a$ есть локальный экстремум и существует $f'(a)$, то $f'(a)=0$.
\textbf{Достаточное условие:} если $f$ непрерывна около $a$, дифференцируема вне $a$ и $f'>0$ слева, $f'<0$ справа от $a$, то $a$ --- строгий локальный максимум. При обратных знаках --- строгий минимум.
\idea{В точке экстремума односторонние разностные отношения имеют противоположные знаки; смена знака производной управляет сравнением значений по обе стороны.}
\uses{определение производной; теорему Лагранжа}
\begin{enumerate}
\item \textbf{Необходимость для максимума.} При малом $h$ имеем $f(a+h)-f(a)\le0$. Поэтому при $h>0$ разностное отношение не больше нуля, а при $h<0$ --- не меньше нуля.
\item Если двусторонняя производная существует, оба предела равны одному $f'(a)$. Из первого знака следует $f'(a)\le0$, из второго $f'(a)\ge0$. Значит, $f'(a)=0$. Для минимума меняются оба знака, а вывод тот же.
\item \textbf{Достаточность для максимума.} Если $x<a$ близко к $a$, по Лагранжу на $[x,a]$
\[
 f(a)-f(x)=f'(\xi)(a-x)>0,\qquad x<\xi<a.
\]
Если $x>a$, по Лагранжу на $[a,x]$
\[
 f(x)-f(a)=f'(\eta)(x-a)<0,\qquad a<\eta<x.
\]
С обеих сторон $f(x)<f(a)$: максимум строгий. Для минимума рассуждение с обратными знаками аналогично.
\end{enumerate}\finish
\remember{Равенство $f'(a)=0$ необходимо, но недостаточно: у $x^3$ в нуле экстремума нет. При достаточном условии производная в самой точке $a$ не требуется.}

\topic{7}{Теорема Кантора}{§ 1.7, с. 19--20}
\program{Теорема о равномерной непрерывности функции, непрерывной на компакте.}
\claim Непрерывная $f:K\to\R$ на компакте $K\subset\R^m$ равномерно непрерывна:
\[
 \forall\varepsilon>0\ \exists\delta>0\ \forall x,y\in K:
 \|x-y\|<\delta\ \Longrightarrow\ |f(x)-f(y)|<\varepsilon.
\]
\idea{Если общей $\delta$ нет, найдутся близкие пары точек с далёкими значениями. Компактность приводит их к одной предельной точке.}
\uses{последовательную компактность; непрерывность}
\begin{enumerate}
\item Предположим противное. Тогда существует $\varepsilon_0>0$ такое, что для каждого $n$ можно выбрать $x_n,y_n\in K$ с
\[
 \|x_n-y_n\|<1/n,\qquad |f(x_n)-f(y_n)|\ge\varepsilon_0.
\]
\item Выделим $x_{n_k}\to a\in K$ по компактности.
\item По неравенству треугольника
\[
 \|y_{n_k}-a\|\le\|y_{n_k}-x_{n_k}\|+\|x_{n_k}-a\|\to0.
\]
Следовательно, $y_{n_k}\to a$ также.
\item По непрерывности обе последовательности значений стремятся к $f(a)$. Их разность стремится к нулю, что противоречит оценке снизу через $\varepsilon_0$.
\end{enumerate}\finish
\remember{Компактность обеспечивает предельную точку внутри $K$, где можно применить непрерывность.}

\topic{8}{Непрерывные частные производные}{§ 1.8, с. 20--22}
\program{Достаточные условия дифференцируемости функции нескольких переменных.}
\claim Если все $f_{x_j}$ существуют около $a\in\R^m$ и непрерывны в $a$, то $f$ дифференцируема в $a$:
\[
 f(a+h)-f(a)=\sum_{j=1}^m f_{x_j}(a)h_j+o(\|h\|).
\]
\idea{Менять координаты по одной: каждое приращение станет одномерным.}
\uses{теорему Лагранжа; неравенство Коши--Буняковского}
\begin{enumerate}
\item Для малого $h$ положим $a^{(j)}=a+(h_1,\ldots,h_j,0,\ldots,0)$ и $a^{(0)}=a$. Тогда
\[
 f(a+h)-f(a)=\sum_{j=1}^m[f(a^{(j)})-f(a^{(j-1)})].
\]
\item В $j$-й разности меняется лишь $j$-я координата. По Лагранжу разность равна $f_{x_j}(c_j)h_j$, где $c_j$ лежит на соответствующем отрезке. При $h_j=0$ можно взять $c_j=a^{(j)}$.
\item Так как $\|c_j-a\|\le\|h\|$, по непрерывности $\alpha_j(h):=f_{x_j}(c_j)-f_{x_j}(a)\to0$.
\item Вычитая линейную часть, для остатка $r(h)$ получаем
\[
 |r(h)|=\left|\sum_j\alpha_j(h)h_j\right|
 \le\left(\sum_j|\alpha_j(h)|^2\right)^{1/2}\|h\|=o(\|h\|).
\]
Это требуемое определение дифференцируемости.
\end{enumerate}\finish
\remember{Одного существования частных производных недостаточно: в шаге 3 использована непрерывность в $a$.}

\topic{9}{Неявная функция одного уравнения}{§ 1.9, с. 22--23}
\program{Теорема о неявной функции, заданной одним уравнением.}
\claim Если $F\in C^1$ около $(a,b)\in\R^m\times\R$, $F(a,b)=0$, $F_y(a,b)\ne0$, то в некоторых окрестностях $W$ точки $a$ и $V$ точки $b$ уравнение $F(x,y)=0$ равносильно $y=\varphi(x)$, где $\varphi\in C^1(W)$ и
\[
 \varphi_{x_j}(x)=-F_{x_j}(x,\varphi(x))/F_y(x,\varphi(x)).
\]
\idea{К уравнению добавить координаты $x$ и обратить полученное отображение.}
\uses{теорему об обратной функции; цепное правило}
\begin{enumerate}
\item Для $\Phi(x,y)=(x,F(x,y))$ имеем
\[
 D\Phi=\begin{pmatrix}I_m&0\\F_x&F_y\end{pmatrix},\qquad
 \det D\Phi(a,b)=F_y(a,b)\ne0.
\]
По теореме об обратной функции $\Phi$ имеет локальное $C^1$-обратное $\Psi$ около $(a,0)$.
\item Первые $m$ координат $\Phi$ не меняет, поэтому $\Psi(u,v)=(u,g(u,v))$. Положим $\varphi(x)=g(x,0)$: тогда $F(x,\varphi(x))=0$ и $\varphi(a)=b$.
\item Выбираем $W\times V$ внутри области взаимной обратимости. По непрерывности, уменьшая $W$, обеспечиваем $(x,0)$ в области $\Psi$ и $\varphi(x)\in V$ для всех $x\in W$.
\item Если $F(x,y)=0$ в $W\times V$, то $\Phi(x,y)=(x,0)$; единственность обратного образа даёт $y=\varphi(x)$. Обратное уже получено в шаге 2.
\item Дифференцируя $F(x,\varphi(x))=0$, получаем $F_{x_j}+F_y\varphi_{x_j}=0$. По непрерывности $F_y$ не обращается в нуль в достаточно малой окрестности, и деление допустимо.
\end{enumerate}\finish
\remember{Здесь доказывается теорема о неявной функции, а теорема об обратной функции принимается известной.}

\topic{10}{Необходимые условия экстремума нескольких переменных}{§ 1.10, с. 23--24}
\program{Экстремумы функций нескольких переменных. Необходимые условия, достаточные условия.}
\claim Пусть $a$ --- внутренняя точка области определения $f$.
Если $f$ дифференцируема в точке локального экстремума $a$, то $\nabla f(a)=0$.
Если дополнительно $f\in C^2$ около $a$, то для локального минимума гессиан $H=f''(a)$ неотрицательно определён, а для максимума --- неположительно определён:
\[
 h^THh\ge0\quad\text{для всех }h\quad(\text{минимум}).
\]
\idea{Ограничить функцию на любую прямую через точку экстремума и применить одномерные необходимые условия.}
\uses{теорему Ферма; цепное правило; одномерную формулу Тейлора}
\begin{enumerate}
\item Зафиксируем произвольное $h\in\R^m$ и положим $g(t)=f(a+th)$ для малых $t$. Так как $a$ внутренняя и является точкой локального экстремума, $t=0$ --- точка того же экстремума $g$.
\item По Ферма и цепному правилу
\[
 0=g'(0)=\nabla f(a)\cdot h.
\]
Это верно при любом $h$. В частности, подставляя каждый координатный вектор, получаем все $f_{x_j}(a)=0$.
\item Пусть теперь $f\in C^2$ и в $a$ минимум. Тогда $g''(0)=h^THh$. По Тейлору и $g'(0)=0$
\[
 0\le g(t)-g(0)=\tfrac12t^2h^THh+o(t^2).
\]
При $t\ne0$ делим на $t^2>0$ и устремляем $t$ к нулю: получаем $h^THh\ge0$.
\item Для максимума $g(t)-g(0)\le0$, поэтому тем же способом $h^THh\le0$ для всех $h$.
\end{enumerate}\finish
\remember{Это необходимые условия, а не достаточные. Например, $f(x,y)=x^4-y^4$ в нуле имеет нулевые градиент и гессиан, но не имеет экстремума. Для достаточного признака нужны строгая определённость гессиана и стационарность.}

\topic{11}{Формула Ньютона--Лейбница}{§ 1.11.2--1.11.3, с. 26--27}
\program{Свойства интеграла с переменным верхним пределом (непрерывность, дифференцируемость). Формула Ньютона-Лейбница.}
\claim Если $f$ непрерывна на $[a,b]$, а $G$ непрерывна на $[a,b]$ и $G'=f$ внутри, то
\[
 \int_a^b f(t)\,dt=G(b)-G(a).
\]
\idea{Интеграл с верхним пределом сам является первообразной; две первообразные отличаются константой.}
\uses{аддитивность и оценку интеграла; теорему Лагранжа}
\begin{enumerate}
\item Положим $F(x)=\int_a^xf(t)\,dt$. Непрерывная функция ограничена: $|f|\le M$ на отрезке. Поэтому
\[
 |F(y)-F(x)|=\left|\int_x^y f(t)\,dt\right|\le M|y-x|,
\]
и $F$ непрерывна на всём отрезке.
\item Для внутренней точки $x$ и малого $h\ne0$
\[
 \frac{F(x+h)-F(x)}h-f(x)
 =\frac1h\int_x^{x+h}[f(t)-f(x)]\,dt.
\]
Модуль правой части не больше максимального отклонения $|f(t)-f(x)|$ на коротком отрезке между $x$ и $x+h$. По непрерывности это отклонение стремится к нулю. Значит, $F'(x)=f(x)$.
\item У $H=G-F$ производная внутри отрезка равна нулю. По Лагранжу для любых $u<v$ из $[a,b]$ имеем $H(v)-H(u)=H'(\xi)(v-u)=0$. Значит, $H$ постоянна.
\item Поэтому $G(b)-F(b)=G(a)-F(a)$. Но $F(a)=0$, а $F(b)=\int_a^bf$. Перегруппировка даёт утверждение.
\end{enumerate}\finish
\remember{Здесь доказаны и непрерывность интеграла с верхним пределом, и его дифференцируемость для непрерывной $f$; затем из них получена сама формула.}

\topic{12}{Почленное интегрирование функционального ряда}{§ 1.12.3, с. 30--32}
\program{Равномерная сходимость функциональных последовательностей и рядов. Непрерывность, интегрируемость и дифференцируемость суммы функционального ряда.}
\claim Если все $u_n$ интегрируемы по Риману на $[a,b]$, а $\sum_{n=1}^\infty u_n(x)$ сходится на этом отрезке равномерно к $S(x)$, то $S$ интегрируема и
\[
 \int_a^b S(x)\,dx=\sum_{n=1}^\infty\int_a^b u_n(x)\,dx.
\]
\idea{Равномерная ошибка по значениям даёт малую ошибку верхних и нижних сумм и малую ошибку интеграла.}
\uses{критерий Дарбу интегрируемости; оценку интеграла}
\begin{enumerate}
\item Обозначим $S_N=\sum_{n=1}^Nu_n$ и $\delta_N=\sup_{[a,b]}|S-S_N|\to0$. Каждая $S_N$ интегрируема и ограничена. Поэтому и $S$ ограничена (берём любой $N$ с конечной $\delta_N$).
\item Для любого разбиения $T$ оценки $S_N-\delta_N\le S\le S_N+\delta_N$ дают
\[
 U(S,T)-L(S,T)\le U(S_N,T)-L(S_N,T)+2\delta_N(b-a),
\]
где $U,L$ --- верхняя и нижняя суммы Дарбу.
\item При заданном $\varepsilon>0$ сначала выберем $N$ так, чтобы $2\delta_N(b-a)<\varepsilon/2$, затем разбиение $T$ с разностью сумм для $S_N$ меньше $\varepsilon/2$. Получаем разность сумм для $S$ меньше $\varepsilon$. По Дарбу $S$ интегрируема. При $a=b$ всё очевидно.
\item Теперь можно применить оценку интеграла:
\[
 \left|\int_a^bS-\int_a^bS_N\right|\le(b-a)\delta_N\to0.
\]
Так как $\int S_N=\sum_{n=1}^N\int u_n$, переход к пределу даёт нужную формулу.
\end{enumerate}\finish
\remember{Не только формально меняем местами сумму и интеграл: сначала проверяем интегрируемость предельной функции.}

\topic{13}{Бесконечная дифференцируемость суммы степенного ряда}{§ 1.13.1, с. 33--35}
\program{Степенные ряды. Радиус сходимости. Бесконечная дифференцируемость суммы степенного ряда. Ряд Тейлора.}
\claim Сумма степенного ряда $f(x)=\sum_{n=0}^\infty a_n(x-x_0)^n$ радиуса $R>0$ бесконечно дифференцируема при $|x-x_0|<R$. Для каждого $k\ge1$
\[
 f^{(k)}(x)=\sum_{n=k}^\infty\frac{n!}{(n-k)!}a_n(x-x_0)^{n-k}.
\]
Ряды производных сходятся равномерно на каждом меньшем замкнутом отрезке.
\idea{Любая фиксированная степень $n$ в коэффициенте производной подавляется геометрическим множителем $q^n$ с $q<1$.}
\uses{признак Вейерштрасса; теорему о почленном дифференцировании (сходимость в одной точке и равномерная сходимость ряда производных)}
\begin{enumerate}
\item Выберем $0<r<\rho<R$. Ряд сходится в $x_0+\rho$, поэтому его члены ограничены: $|a_n|\rho^n\le M$ для всех $n$.
\item Для $|x-x_0|\le r$ и $n\ge k$ модуль $k$-й производной члена не больше
\[
 \frac{n!}{(n-k)!}|a_n|r^{n-k}
 \le Mr^{-k}n^k(r/\rho)^n.
\]
Здесь $n!/(n-k)!\le n^k$.
\item При $q=r/\rho<1$ числовой ряд $\sum n^kq^n$ сходится: отношение соседних членов стремится к $q<1$. По Вейерштрассу каждый формальный ряд производных сходится равномерно на этом отрезке.
\item Применяем теорему о почленном дифференцировании сначала к исходному ряду, затем к ряду первых производных и так далее. В центре каждый из этих рядов сходится: при подстановке $x=x_0$ остаётся лишь один ненулевой член. Условия теоремы на каждом шаге выполнены.
\item Получаем все производные и указанные формулы. Произвольная внутренняя точка лежит на некотором отрезке $|x-x_0|\le r<R$, поэтому результат верен внутри всего интервала сходимости.
\end{enumerate}\finish
\remember{В частности, $f^{(k)}(x_0)=k!a_k$: исходный степенной ряд является рядом Тейлора своей суммы. На границе радиуса это рассуждение не действует.}

\topic{14}{Формула Грина}{§ 1.16, с. 43--44}
\program{Формула Грина. Потенциальные векторные поля на плоскости.}
\claim Если $D$ --- ограниченная плоская область с кусочно-гладкой границей (со стандартными регулярными углами), $P,Q\in C^1$ в окрестности $\overline D$, то
\[
 \oint_{\partial D}P\,dx+Q\,dy
 =\iint_D(Q_x-P_y)\,dx\,dy.
\]
Внешняя граница обходится против часовой стрелки, отверстия --- по часовой.
\idea{Около каждого гладкого участка границы свести формулу к Ньютону--Лейбницу и правилу производной интеграла; затем сложить локальные равенства.}
\uses{Ньютона--Лейбница; Фубини; геометрическую локализацию: гладкое разбиение единицы по внутренним участкам и участкам, где граница после поворота осей является графиком}
\begin{enumerate}
\item \textbf{Локальный случай.} Пусть в выбранном участке область лежит ниже $y=g(x)$, а $P,Q$ равны нулю около всех краёв участка, кроме этого графика. Функция $g$ может быть кусочно $C^1$. Выберем нижний уровень $c$ вне носителей $P,Q$.
\item Сначала интегрируем по $y$:
\[
 \iint_D(-P_y)\,dx\,dy=-\int P(x,g(x))\,dx.
\]
На верхней границе положительный обход идёт справа налево, поэтому правая часть равна $\int_{\partial D}P\,dx$.
\item По правилу производной интеграла
\[
 \frac d{dx}\int_c^{g(x)}Q(x,y)\,dy
 =\int_c^{g(x)}Q_x(x,y)\,dy+Q(x,g(x))g'(x).
\]
Интеграл полной производной нулевой, так как интегрируемая функция равна нулю у краёв участка. Поэтому
\[
 \iint_D Q_x\,dx\,dy=-\int Q(x,g(x))g'(x)\,dx
 =\int_{\partial D}Q\,dy.
\]
Сумма шагов 2--3 даёт локальную формулу. Для области выше графика оба граничных знака меняются согласованно; поворот осей сохраняет формулу.
\item \textbf{Склейка.} Берём гладкие функции $\chi_j$ с носителями в таких участках и внутри $D$, причём $\sum_j\chi_j=1$ около $\overline D$. Применяем локальное равенство к $P_j=\chi_jP$, $Q_j=\chi_jQ$. Для внутреннего носителя обе стороны равны нулю по Ньютону--Лейбницу. Суммируя, используем
\[
 \sum_jP_j=P,\qquad\sum_jQ_j=Q,
\]
и их производные. Получаем формулу на всей области.
\end{enumerate}\finish
\remember{Геометрическая локализация --- отдельный стандартный инструмент, принятый здесь без доказательства. Сама формула Грина выведена из одномерного интегрирования, а не заменена следствием о потенциальности.}

\topic{15}{Формула Остроградского--Гаусса}{§ 1.17, с. 44}
\program{Формула Остроградского- Гаусса. Соленоидальные векторные поля.}
\claim Пусть $V\subset\R^3$ --- ограниченная область с гладкой границей, $F=(P,Q,R)\in C^1$ около $\overline V$. При внешней нормали
\[
 \iint_{\partial V}F\cdot n\,dS
 =\iiint_V(P_x+Q_y+R_z)\,dx\,dy\,dz.
\]
Формула распространяется сложением на области с границей из конечного числа согласованных гладких кусков.
\idea{Для участка границы $z=g(x,y)$ три объёмных интеграла дают три компоненты потока; локальные результаты складываются.}
\uses{Фубини; Ньютона--Лейбница; дифференцирование интеграла; разбиение единицы по внутренним участкам и граничным графикам}
\begin{enumerate}
\item \textbf{Один участок.} Пусть $V$ локально лежит ниже графика $z=g(x,y)$, а поле имеет носитель внутри этого участка. Ниже носителя выберем уровень $c$.
\item По правилу производной интеграла
\[
 \int_c^{g(x,y)}P_x\,dz
 =\partial_x\left(\int_c^{g(x,y)}P\,dz\right)-P(x,y,g)g_x.
\]
Интегрируя по $(x,y)$, первый член получаем нулевым: функция под полной производной исчезает у краёв носителя. Поэтому
\[
 \iiint_V P_x=-\iint P(x,y,g)g_x\,dx\,dy.
\]
Аналогично $\iiint_VQ_y=-\iint Q(x,y,g)g_y\,dx\,dy$.
\item По Ньютону--Лейбницу $\iiint_VR_z=\iint R(x,y,g)\,dx\,dy$. Складывая, получаем
\[
 \iiint_V\operatorname{div}F
 =\iint[-Pg_x-Qg_y+R]_{z=g}\,dx\,dy.
\]
Для области ниже графика внешняя нормаль задаёт вектор площади
\[
 n\,dS=(-g_x,-g_y,1)\,dx\,dy.
\]
Поэтому правая часть --- ровно поток. Для области выше графика знаки меняются согласованно.
\item \textbf{Общий случай.} Повороты координат сохраняют дивергенцию и поток. Разложим $F=\sum_j\chi_jF$ разбиением единицы. Для каждого граничного участка применяем доказанное; для внутреннего носителя интеграл дивергенции равен нулю по Фубини и Ньютону--Лейбницу. Складывая и используя $\sum\chi_j=1$, получаем всю формулу.
\end{enumerate}\finish
\remember{Разбиение единицы --- используемый геометрический факт, а не скрытое применение Гаусса. Для кусочной границы при сложении поток через общую внутреннюю грань сокращается из-за противоположных нормалей.}

\topic{16}{Формула Стокса}{§ 1.15, с. 42--43}
\program{Формула Стокса.}
\claim Пусть $S$ --- компактная ориентированная гладкая поверхность с кусочно-гладким краем, $F\in C^1$ в её окрестности. Тогда при согласованных ориентациях
\[
 \oint_{\partial S}F\cdot dr=\iint_S(\operatorname{rot}F)\cdot n\,dS.
\]
\idea{Перенести криволинейный интеграл на плоскость параметров и применить Грина.}
\uses{формулу Грина (раздел 14); цепное правило; конечное разбиение гладкой поверхности на согласованно ориентированные параметрические куски}
\begin{enumerate}
\item \textbf{Один кусок.} Запишем $S_j$ регулярной $C^2$-параметризацией $r(u,v)$ над плоской областью $D$ с кусочно-гладкой границей. При $dr=r_u\,du+r_v\,dv$ положим $A=F(r)\cdot r_u$, $B=F(r)\cdot r_v$. Тогда по определению интеграла
\[
 \oint_{\partial S_j}F\cdot dr=\oint_{\partial D}A\,du+B\,dv
 =\iint_D(B_u-A_v)\,du\,dv.
\]
Граница $\partial D$ ориентирована положительно.
\item Цепное правило даёт
\[
 B_u-A_v=(DF(r)r_u)\cdot r_v-(DF(r)r_v)\cdot r_u.
\]
Члены $F(r)\cdot r_{vu}$ и $F(r)\cdot r_{uv}$ сократились, поскольку $r\in C^2$.
\item Для любых векторов $p,q$ верно тождество
\[
 (DF\,p)\cdot q-(DF\,q)\cdot p
 =(\operatorname{rot}F)\cdot(p\times q).
\]
Оно проверяется раскрытием трёх координат. С $p=r_u$, $q=r_v$ получаем нужный подынтегральный член.
\item По определению ориентированного поверхностного интеграла $n\,dS=(r_u\times r_v)\,du\,dv$. Подстановка доказывает формулу для одного куска $S_j$.
\item \textbf{Сложение кусков.} Разбиваем поверхность на конечное число таких кусков с согласованными ориентациями и суммируем полученные равенства. Поверхностные интегралы складываются в интеграл по $S$.
\item Каждая внутренняя линия разреза является границей двух соседних кусков. Согласованная ориентация задаёт на ней противоположные направления обхода, поэтому два криволинейных интеграла сокращаются. Остаётся только интеграл по настоящему краю $\partial S$. Это и есть формула для всей поверхности.
\end{enumerate}\finish
\remember{Локальное вычисление и сокращение внутренних разрезов --- две отдельные части доказательства. Возможность конечного разбиения поверхности на гладкие параметрические куски здесь является используемым геометрическим фактом.}

\topic{17}{Признак Дини сходимости ряда Фурье}{§ 1.18.2--1.18.3, с. 45--47}
\program{Достаточные условия сходимости тригонометрического ряда Фурье в точке.}
\claim Пусть $f$ --- $2\pi$-периодическая, $f\in L^1[-\pi,\pi]$. Если для некоторых $S$ и $\delta>0$
\[
 \int_0^\delta\frac{|f(x+t)+f(x-t)-2S|}{t}\,dt<\infty,
\]
то частичные суммы её ряда Фурье удовлетворяют $S_n(f,x)\to S$.
\idea{Представить ошибку как осциллирующий интеграл от одной интегрируемой функции.}
\uses{ядро Дирихле; лемму Римана--Лебега: для $g\in L^1(0,\pi)$ интеграл $\int_0^\pi g(t)\sin(\lambda t)\,dt\to0$ при $\lambda\to\infty$}
\begin{enumerate}
\item В нормировке конспекта
\[
 D_n(t)=\tfrac12+\sum_{k=1}^n\cos kt
 =\frac{\sin((n+1/2)t)}{2\sin(t/2)},\qquad
 \int_0^\pi D_n(t)\,dt=\frac\pi2.
\]
Формула свёртки для частичной суммы даёт
\[
 S_n(f,x)-S=\frac1\pi\int_0^\pi
 g(t)\sin((n+1/2)t)\,dt,
\]
где $g(t)=[f(x+t)+f(x-t)-2S]/[2\sin(t/2)]$.
\item При $0<t\le\pi$ имеем $\sin(t/2)\ge t/\pi$. Поэтому около нуля
\[
 |g(t)|\le\frac\pi2\frac{|f(x+t)+f(x-t)-2S|}{t}.
\]
Правая часть интегрируема по условию (при необходимости уменьшаем $\delta$ до числа меньше $\pi$).
\item Вдали от нуля знаменатель отделён от нуля; числитель интегрируем по периодичности и $f\in L^1$. Итак, $g\in L^1(0,\pi)$.
\item По лемме Римана--Лебега интеграл ошибки стремится к нулю. Значит, $S_n(f,x)\to S$.
\end{enumerate}\finish
\remember{Основная работа --- проверить интегрируемость $g$ около нуля; именно это обеспечивает условие Дини.}

\topic{18}{Равномерная сходимость ряда Фурье при условии Гёльдера}{§ 1.19, с. 47--48}
\program{Достаточные условия равномерной сходимости тригонометрического ряда Фурье.}
\claim Если $2\pi$-периодическая функция $f$ удовлетворяет условию Гёльдера
\[
 |f(x)-f(y)|\le C|x-y|^\alpha,\qquad x,y\in\R,\quad0<\alpha\le1,
\]
то её ряд Фурье сходится к $f$ равномерно на $\R$.
\idea{Ошибка выражается осциллирующим интегралом. Условие Гёльдера даёт общую интегрируемую оценку около нуля; компактность позволяет сделать сходимость интеграла общей для всех $x$.}
\uses{ядро Дирихле; лемму Римана--Лебега; мажорируемую сходимость; компактность непрерывного образа отрезка}
\begin{enumerate}
\item Как в доказательстве Дини, имеем
\[
 S_n(f,x)-f(x)=\frac1\pi\int_0^\pi g_x(t)\sin((n+1/2)t)\,dt,
\]
где $g_x(t)=[f(x+t)+f(x-t)-2f(x)]/[2\sin(t/2)]$.
\item Условие Гёльдера и $\sin(t/2)\ge t/\pi$ дают общую для всех $x$ оценку
\[
 |g_x(t)|\le\pi C t^{\alpha-1},\qquad0<t\le\pi.
\]
Правая часть интегрируема, потому что $\alpha>0$.
\item Если $x_j\to x$, непрерывность $f$ даёт $g_{x_j}(t)\to g_x(t)$ для $t>0$. Разности по модулю не больше $2\pi Ct^{\alpha-1}$. По мажорируемой сходимости
\[
 \|g_{x_j}-g_x\|_1:=\int_0^\pi|g_{x_j}(t)-g_x(t)|\,dt\longrightarrow0.
\]
Значит, отображение $x\mapsto g_x$ непрерывно из $[0,2\pi]$ в $L^1(0,\pi)$.
\item Его образ компактен. Поэтому для заданного $\varepsilon>0$ можно выбрать конечное число точек $x_1,\ldots,x_m$ так, чтобы для каждого $x\in[0,2\pi]$ нашёлся $j$ с $\|g_x-g_{x_j}\|_1<\pi\varepsilon/2$. Это свойство конечного покрытия компактного множества шарами фиксированного радиуса.
\item Для каждой из конечного числа $g_{x_j}$ лемма Римана--Лебега даёт малость осциллирующего интеграла. Выбрав один достаточно большой $N$, при $n\ge N$ имеем одновременно для всех $j$
\[
 \left|\int_0^\pi g_{x_j}(t)\sin((n+1/2)t)\,dt\right|<\pi\varepsilon/2.
\]
\item При произвольном $x$ заменяем $g_x$ на подходящую $g_{x_j}$. Ошибка интеграла не больше $\|g_x-g_{x_j}\|_1$, поскольку $|\sin|\le1$. Две ошибки дают $|S_n(f,x)-f(x)|<\varepsilon$ для всех $x\in[0,2\pi]$. Периодичность распространяет оценку на $\R$.
\end{enumerate}\finish
\remember{Главное отличие от поточечного Дини --- шаги 3--5: они позволяют выбрать один номер $N$ сразу для всех точек. Одной непрерывности без количественного условия Гёльдера для этого вывода недостаточно.}

\topic{19}{Три основных свойства преобразования Фурье}{§ 1.20.2, с. 49--51}
\program{Непрерывность преобразования Фурье абсолютно интегрируемой функции. Преобразование Фурье производной и производная преобразования Фурье.}
\claim В нормировке $\widehat f(y)=(2\pi)^{-1/2}\int_\R f(x)e^{-ixy}\,dx$ верны:
\begin{enumerate}[label=\alph*)]
\item при $f\in L^1(\R)$ функция $\widehat f$ непрерывна;
\item при $f\in C^1(\R)$, $f,f'\in L^1(\R)$ имеем $\widehat{f'}(y)=iy\widehat f(y)$;
\item при $f,xf\in L^1(\R)$ имеем $(\widehat f)'(y)=-i\widehat{xf}(y)$.
\end{enumerate}
\idea{Непрерывность и производная следуют из мажорируемого перехода к пределу; преобразование производной --- из интегрирования по частям.}
\uses{теорему о мажорируемой сходимости; интегрирование по частям}
\begin{enumerate}
\item \textbf{Непрерывность.} При $h\to0$ функция $f(x)e^{-ix(y+h)}$ поточечно стремится к $f(x)e^{-ixy}$ и по модулю не больше $|f(x)|\in L^1$. Поэтому можно перейти к пределу под интегралом: $\widehat f(y+h)\to\widehat f(y)$.
\item \textbf{Сначала проверим граничные значения.} Из $f'\in L^1$ и равенства $f(x)=f(0)+\int_0^xf'$ следует существование конечных пределов $f(\pm\infty)$. Из $f\in L^1$ эти пределы равны нулю: ненулевой предел дал бы $|f(x)|\ge c>0$ на бесконечном хвосте.
\item Интегрирование по частям на $[-A,A]$ даёт
\[
 \int_{-A}^A f'(x)e^{-ixy}\,dx
 =[f(x)e^{-ixy}]_{-A}^A+iy\int_{-A}^Af(x)e^{-ixy}\,dx.
\]
При $A\to\infty$ граничный член исчезает по шагу 2, оба интеграла сходятся абсолютно. Делим на $\sqrt{2\pi}$ и получаем пункт б).
\item \textbf{Производная преобразования.} Для $h\ne0$
\[
 \frac{\widehat f(y+h)-\widehat f(y)}h
 =\frac1{\sqrt{2\pi}}\int_\R f(x)e^{-ixy}\frac{e^{-ixh}-1}{h}\,dx.
\]
При $h\to0$ дробь стремится к $-ix$. Неравенство $|e^{it}-1|\le|t|$ даёт интегрируемую мажоранту $|xf(x)|$. Переходя к пределу под интегралом, получаем пункт в).
\end{enumerate}\finish
\remember{Для пункта б) нужна интегрируемость $f'$, для пункта в) --- интегрируемость $xf$. Эти условия отвечают разным операциям.}

\topic{20}{Расстояние от точки до плоскости}{§ 2.1.3, с. 53--54}
\program{Прямые и плоскости в пространстве. Формулы расстояния от точки до прямой и плоскости, между прямыми в пространстве. Углы между прямыми и плоскостями.}
\claim Для плоскости $\Pi:Ax+By+Cz+D=0$, где $n=(A,B,C)\ne0$, и точки $M=(x_0,y_0,z_0)$
\[
 d(M,\Pi)=\frac{|Ax_0+By_0+Cz_0+D|}{\sqrt{A^2+B^2+C^2}}.
\]
\idea{Найти основание перпендикуляра вдоль нормали и доказать, что оно ближе всех остальных точек плоскости.}
\uses{скалярное произведение; теорему Пифагора}
\begin{enumerate}
\item Ищем $H=M-tn$. Условие $H\in\Pi$ равносильно
\[
 n\cdot M-t\|n\|^2+D=0,
 \quad t=\frac{n\cdot M+D}{\|n\|^2}.
\]
Деление допустимо, поскольку $n\ne0$.
\item Вектор $M-H=tn$ перпендикулярен плоскости. Для любой $X\in\Pi$ имеем $n\cdot(X-H)=0$, так как обе точки удовлетворяют уравнению плоскости.
\item Поэтому $M-H$ и $H-X$ ортогональны, и
\[
 \|M-X\|^2=\|M-H\|^2+\|H-X\|^2\ge\|M-H\|^2.
\]
Минимум расстояния достигается именно при $X=H$.
\item Наконец,
\[
 d(M,\Pi)=\|M-H\|=|t|\|n\|
 =\frac{|n\cdot M+D|}{\|n\|},
\]
что и даёт формулу в координатах.
\end{enumerate}\finish
\remember{Недостаточно найти длину нормальной проекции: шаг 3 объясняет, почему это действительно наименьшее расстояние.}

\topic{21}{Эллипс как множество точек с постоянной суммой расстояний}{§ 2.2.2, с. 56--57}
\program{Кривые второго порядка, их геометрические свойства.}
\claim Пусть $a>b>0$, $c=\sqrt{a^2-b^2}$, $F_1=(-c,0)$, $F_2=(c,0)$. Тогда
\[
 \frac{x^2}{a^2}+\frac{y^2}{b^2}=1
 \quad\Longleftrightarrow\quad |MF_1|+|MF_2|=2a,
 \qquad M=(x,y).
\]
\idea{В прямую сторону квадраты расстояний превращаются в полные квадраты; в обратную --- используем сумму и разность расстояний.}
\uses{формулу расстояния; $a^2=b^2+c^2$}
\begin{enumerate}
\item \textbf{Из канонического уравнения.} Пусть $x^2/a^2+y^2/b^2=1$. Тогда $|x|\le a$ и $y^2=b^2(1-x^2/a^2)$. Поэтому
\[
 r_1^2=(x+c)^2+y^2=(a+cx/a)^2,
 \quad r_2^2=(x-c)^2+y^2=(a-cx/a)^2.
\]
Так как $a\pm cx/a\ge a-c>0$, имеем $r_1=a+cx/a$, $r_2=a-cx/a$, и их сумма равна $2a$.
\item \textbf{Из фокального условия.} Пусть $r_1+r_2=2a$. Из формул расстояния $r_1^2-r_2^2=4cx$. Делим на ненулевую сумму:
\[
 r_1-r_2=2cx/a.
\]
Складывая это равенство с $r_1+r_2=2a$, получаем $r_1=a+cx/a$.
\item Подставляем в $r_1^2=(x+c)^2+y^2$:
\[
 a^2+2cx+c^2x^2/a^2=x^2+2cx+c^2+y^2.
\]
После сокращения и использования $a^2-c^2=b^2$ имеем $b^2=(b^2/a^2)x^2+y^2$. Деление на $b^2>0$ даёт уравнение эллипса.
\end{enumerate}\finish
\remember{Доказаны обе стороны: геометрическое определение эллипса и его каноническое уравнение действительно описывают одно множество.}

\topic{22}{Теорема Кронекера--Капелли}{§ 2.3.2, с. 57--58}
\program{Общее решение системы линейных алгебраических уравнений. Теорема Кронекера- Капелли.}
\claim Система $Ax=b$, где $A$ --- матрица размера $m\times n$, совместна тогда и только тогда, когда $\operatorname{rank}A=\operatorname{rank}(A\mid b)$.
\idea{Решение существует ровно тогда, когда добавленный столбец уже выражается через старые и не увеличивает их линейную оболочку.}
\uses{ранг как размерность пространства столбцов; свойство размерности вложенных подпространств}
\begin{enumerate}
\item Обозначим столбцы $A$ через $a_1,\ldots,a_n\in\R^m$. Матричное уравнение $Ax=b$ означает
\[
 x_1a_1+\cdots+x_na_n=b.
\]
Следовательно, система совместна ровно при $b\in L$, где $L=\operatorname{span}(a_1,\ldots,a_n)$.
\item Пространство столбцов расширенной матрицы равно $L'=\operatorname{span}(L,b)$; всегда $L\subset L'$.
\item Если решение существует, то $b\in L$. Поэтому $L'=L$ и ранги совпадают.
\item Обратно, равенство рангов означает $\dim L=\dim L'$. Для конечномерных пространств включение и равенство размерностей дают $L=L'$: иначе добавление вектора из $L'\setminus L$ к базису $L$ увеличило бы размерность.
\item Значит, $b\in L$. По определению линейной оболочки существуют коэффициенты $x_1,\ldots,x_n$, выражающие $b$ через столбцы $A$. Это и есть решение системы.
\end{enumerate}\finish
\remember{Равенство рангов даёт существование решения. Единственность --- отдельный вопрос: для неё нужен ещё полный столбцовый ранг $n$.}

\topic{23}{Теорема о размерности ядра и образа}{§ 2.4.2, с. 59--60}
\program{Линейное пространство, базис и размерность. Линейное отображение конечномерных пространств, его матрица. Ядро и образ линейного отображения.}
\claim Для линейного $T:U\to V$ с конечномерным $U$
\[
 \dim U=\dim\ker T+\dim\operatorname{Im}T.
\]
\idea{Дополнить базис ядра до базиса всего пространства; образы добавленных векторов образуют базис образа.}
\uses{дополнение линейно независимой системы до базиса}
\begin{enumerate}
\item Пусть $e_1,\ldots,e_k$ --- базис $\ker T$. Дополним его векторами $e_{k+1},\ldots,e_n$ до базиса $U$. Здесь $n=\dim U$.
\item Любой $u\in U$ имеет вид $\sum_{j=1}^n\alpha_je_j$. Так как $Te_j=0$ при $j\le k$, имеем
\[
 Tu=\sum_{j=k+1}^n\alpha_jTe_j.
\]
Значит, $Te_{k+1},\ldots,Te_n$ порождают образ.
\item Докажем их независимость. Если $\sum_{j=k+1}^n\beta_jTe_j=0$, то $w=\sum_{j=k+1}^n\beta_je_j\in\ker T$. Поэтому также $w=\sum_{j=1}^k\gamma_je_j$.
\item Вычитая два выражения $w$, получаем линейную комбинацию полного базиса, равную нулю. По независимости базиса все $\beta_j$ равны нулю. Независимость образов доказана.
\item Таким образом, образ имеет базис из $n-k$ векторов. Следовательно, $\dim\operatorname{Im}T=n-k$, откуда формула.
\end{enumerate}\finish
\remember{Самая важная часть --- независимость образов. Она неверна для произвольно выбранных векторов: здесь помогает дополнение именно базиса ядра.}

\topic{24}{Критерий диагонализируемости}{§ 2.5, с. 60--61}
\program{Собственные значения и собственные векторы линейных преобразований. Диагонализируемость линейных преобразований.}
\claim Линейный оператор $T$ в $n$-мерном пространстве над $\R$ или $\C$ диагонализируем тогда и только тогда, когда характеристический многочлен раскладывается над этим полем на линейные множители и для каждого собственного значения
\[
 \dim\ker(T-\lambda I)=m_\lambda,
\]
где $m_\lambda$ --- кратность корня $\lambda$ в характеристическом многочлене.
\idea{Диагональная матрица сразу даёт нужные размерности; обратно, базисы собственных подпространств объединяются в базис всего пространства.}
\uses{характеристический многочлен; независимость собственных подпространств для разных значений}
\begin{enumerate}
\item \textbf{Необходимость.} В базисе собственных векторов матрица диагональна. Характеристический многочлен есть произведение $(t-\lambda_j)$ по всем диагональным элементам. Число появлений значения $\lambda$ равно и его кратности, и числу свободных координат решения $(T-\lambda I)v=0$. Получаем требуемое равенство.
\item \textbf{Почему собственные подпространства независимы.} Если $v_j\in\ker(T-\lambda_jI)$ и $\sum_{j=1}^sv_j=0$ при разных $\lambda_j$, применим для фиксированного $k$ оператор $\prod_{j\ne k}(T-\lambda_jI)$. Он уничтожает все члены, кроме $k$-го, и даёт
\[
 \left(\prod_{j\ne k}(\lambda_k-\lambda_j)\right)v_k=0.
\]
Множитель ненулевой, поэтому $v_k=0$. Это верно для каждого $k$.
\item \textbf{Достаточность.} Выбираем базис в каждом собственном подпространстве. По шагу 2 их объединение линейно независимо. Векторов в нём
\[
 \sum_\lambda\dim\ker(T-\lambda I)=\sum_\lambda m_\lambda=n,
\]
поскольку характеристический многочлен полностью разложился. Значит, это базис всего пространства из собственных векторов. В нём матрица диагональна.
\end{enumerate}\finish
\remember{Повторяющиеся собственные значения допускаются. Важно равенство геометрической и алгебраической кратностей, а также разложимость над выбранным полем.}

\topic{25}{Спектральная теорема для самосопряжённого оператора}{§ 2.6, с. 62--63}
\program{Самосопряженные преобразования евклидовых пространств, свойства их собственных значений и собственных векторов.}
\claim В конечномерном вещественном евклидовом пространстве у самосопряжённого $T$ существует ортонормированный базис собственных векторов. Поэтому его матрица в таком базисе диагональна и все собственные значения вещественны.
\idea{Найти один собственный вектор максимизацией $(Tx,x)$ на сфере, затем повторить на его ортогональном дополнении.}
\uses{Вейерштрасса на сфере; Ферма; индукцию по размерности}
\begin{enumerate}
\item При размерности 1 всё очевидно. В общем случае непрерывная вещественная функция $q(x)=(Tx,x)$ достигает максимума на единичной сфере в некотором $v$, $\|v\|=1$.
\item Возьмём любой единичный $w\perp v$. Кривая $x(t)=v\cos t+w\sin t$ лежит на сфере. У $q(x(t))$ максимум при $t=0$, поэтому
\[
 0=\left.\frac d{dt}(Tx(t),x(t))\right|_{t=0}
 =(Tw,v)+(Tv,w)=2(Tv,w).
\]
Последнее равенство --- самосопряжённость. Значит, $Tv$ ортогонален всему $v^\perp$, то есть $Tv=\lambda v$. Число $\lambda=(Tv,v)$ вещественно.
\item Подпространство $v^\perp$ инвариантно: для $u\perp v$
\[
 (Tu,v)=(u,Tv)=\lambda(u,v)=0.
\]
Ограничение $T$ на $v^\perp$ остаётся самосопряжённым, поскольку равенство скалярных произведений сохраняется.
\item По индукции в $(n-1)$-мерном $v^\perp$ есть ортонормированный базис собственных векторов ограничения. Добавляем $v$ и получаем нужный базис всего пространства.
\end{enumerate}\finish
\remember{Для вещественной симметричной матрицы это означает $Q^TAQ=\operatorname{diag}(\lambda_1,\ldots,\lambda_n)$ с ортогональной $Q$. Это существенно сильнее одного свойства ортогональности разных собственных векторов.}

\topic{26}{Критерий Сильвестра}{§ 2.7.2, с. 64--65}
\program{Приведение квадратичных форм в линейном пространстве к каноническому виду. Положительно определенные квадратичные формы. Критерий Сильвестра.}
\claim Для вещественной симметричной $n\times n$ матрицы $A$ форма $q(x)=x^TAx$ положительно определена тогда и только тогда, когда все угловые главные миноры $\Delta_1,\ldots,\Delta_n$ положительны.
\idea{Выделить полный квадрат по первой переменной. Оставшаяся форма имеет миноры, выражающиеся через исходные.}
\uses{выделение полного квадрата; блочную формулу определителя; индукцию}
\begin{enumerate}
\item При $n=1$ утверждение очевидно. Пусть оно доказано для размера $n-1$. Запишем
\[
 A=\begin{pmatrix}a&b^T\\b&C\end{pmatrix},\qquad x=(s,z).
\]
При $a>0$ выделение квадрата даёт
\[
 q(s,z)=a\left(s+\frac{b^Tz}{a}\right)^2+z^TBz,
 \quad B=C-\frac{bb^T}{a}.
\]
\item Матрица $B$ симметрична. Для её углового минора размера $j$ блочная формула (или вычитание кратных первой строки) даёт
\[
 \Delta_{j+1}(A)=a\Delta_j(B),\qquad1\le j\le n-1.
\]
\item \textbf{Достаточность.} Если все $\Delta_j(A)>0$, то $a=\Delta_1>0$ и все $\Delta_j(B)>0$. По индукции $z^TBz>0$ при $z\ne0$. В формуле шага 1 при $z\ne0$ положителен второй член, а при $z=0$, $s\ne0$ --- первый. Итак, $q>0$ вне нуля.
\item \textbf{Необходимость.} Если $q$ положительно определена, то $a=q(1,0)>0$. Для $z\ne0$ подставим $s=-b^Tz/a$: получим $z^TBz=q(s,z)>0$. По индукции все миноры $B$ положительны. Формула шага 2 даёт все $\Delta_j(A)>0$.
\end{enumerate}\finish
\remember{Речь именно об угловых главных минорах. Замена строгого знака на $\ge0$ в этом критерии не даёт критерия полуопределённости.}

\topic{27}{Частное решение для правой части-квазимногочлена}{§ 3.2, с. 71--72}
\program{Линейные обыкновенные дифференциальные уравнения с постоянными коэффициентами и правой частью - квазимногочленом.}
\claim Пусть $M(D)y=e^{\gamma x}P_m(x)$, где $M$ --- ненулевой многочлен с постоянными коэффициентами, $D=d/dx$, $\deg P_m\le m$. Если $\gamma$ --- корень $M$ кратности $s$ (при $M(\gamma)\ne0$ полагаем $s=0$), то существует частное решение
\[
 y_p=e^{\gamma x}x^sQ_m(x),\qquad\deg Q_m\le m.
\]
\idea{После удаления экспоненты задача становится конечной системой на коэффициенты многочлена; её матрица треугольна с ненулевой диагональю.}
\uses{тождество сдвига $M(D)(e^{\gamma x}u)=e^{\gamma x}M(D+\gamma)u$; обратимость треугольной матрицы с ненулевой диагональю}
\begin{enumerate}
\item По кратности корня $M(z+\gamma)=z^sH(z)$, где $H(0)\ne0$. После подстановки $y=e^{\gamma x}u$ уравнение равносильно
\[
 H(D)D^su=P_m.
\]
Множители коммутируют: они многочлены от одного оператора $D$.
\item Ищем $u=x^sQ_m$. Рассмотрим линейное отображение
\[
 L:Q\longmapsto H(D)D^s(x^sQ)
\]
из пространства многочленов степени не больше $m$ в него же.
\item На базисном многочлене $Q=x^k$, $0\le k\le m$, имеем
\[
 D^s(x^{s+k})=\frac{(s+k)!}{k!}x^k.
\]
Оператор $H(D)=H(0)+h_1D+\cdots$ даёт старший член $H(0)(s+k)!x^k/k!$; остальные члены имеют меньшую степень.
\item Значит, матрица $L$ в базисе $1,x,\ldots,x^m$ треугольна, а диагональные элементы $H(0)(s+k)!/k!$ ненулевые. Матрица обратима, поэтому для данного $P_m$ существует $Q_m$ с $LQ_m=P_m$. Возвращаясь к $y$, получаем требуемый вид.
\end{enumerate}\finish
\remember{Множитель $x^s$ --- поправка на резонанс. Для вещественных синусно-косинусных правых частей берут вещественную или мнимую часть комплексного решения; линейность сохраняет равенство.}

\topic{28}{Решение системы через матричную экспоненту}{§ 3.3.2, с. 74--76}
\program{Системы линейных однородных дифференциальных уравнений с постоянными коэффициентами, методы их решения.}
\claim Для постоянной вещественной матрицы $A$ единственное решение задачи $x'=Ax$, $x(t_0)=x_0$ имеет вид
\[
 x(t)=e^{(t-t_0)A}x_0,\qquad e^{tA}=\sum_{k=0}^\infty\frac{t^kA^k}{k!}.
\]
\idea{Ряд экспоненты можно дифференцировать; для единственности умножаем любое решение на обратную экспоненту.}
\uses{признак Вейерштрасса; почленное дифференцирование равномерно сходящихся рядов производных}
\begin{enumerate}
\item На $|t|\le R$ нормы членов ряда не больше $(R\|A\|)^k/k!$. Ряд этих чисел сходится. Для ряда производных мажоранта $\|A\|\sum_{j=0}^\infty(R\|A\|)^j/j!$ тоже сходится.
\item Поэтому допустимо дифференцирование:
\[
 \frac d{dt}e^{tA}=Ae^{tA}=e^{tA}A,\qquad e^{0A}=I.
\]
Коммутирование с $A$ следует из того, что каждый член ряда есть степень $A$.
\item Производная $e^{-tA}e^{tA}$ равна нулю по правилу произведения. При $t=0$ произведение равно $I$, значит, $e^{-tA}$ обратна к $e^{tA}$.
\item Предложенная функция удовлетворяет уравнению и начальному условию по шагу 2. Если $z$ --- любое другое решение, то
\[
 \frac d{dt}\bigl(e^{-(t-t_0)A}z(t)\bigr)
 =e^{-(t-t_0)A}(-Az+z')=0.
\]
Следовательно, произведение постоянно и равно $x_0$. Умножаем на $e^{(t-t_0)A}$ и получаем ту же формулу.
\end{enumerate}\finish
\remember{Сходимости самого ряда недостаточно для дифференцирования: в шаге 1 отдельно проверяется ряд производных.}

\topic{29}{Формула Лиувилля--Остроградского}{§ 3.1, с. 66--68}
\program{Линейные обыкновенные дифференциальные уравнения с переменными коэффициентами. Фундаментальная система решений. Определитель Вронского. Формула Лиувилля-Остроградского.}
\claim Для решений $y_1,\ldots,y_n$ уравнения
\[
 y^{(n)}+p_1(t)y^{(n-1)}+\cdots+p_n(t)y=0
\]
с непрерывными коэффициентами на интервале $I$ их вронскиан удовлетворяет
\[
 W(t)=W(t_0)\exp\left(-\int_{t_0}^t p_1(s)\,ds\right).
\]
\idea{При дифференцировании определителя все слагаемые, кроме последнего, имеют одинаковые строки.}
\uses{правило производной определителя; линейность по строке}
\begin{enumerate}
\item Обозначим $R_j=(y_1^{(j)},\ldots,y_n^{(j)})$. Тогда $W=\det(R_0,\ldots,R_{n-1})$, где перечислены строки.
\item Производная определителя --- сумма определителей, в каждом из которых дифференцируется одна строка. Если дифференцируется $R_j$, $j<n-1$, вместо неё возникает уже имеющаяся строка $R_{j+1}$. Такой определитель нулевой.
\item Остаётся производная последней строки. По уравнению
\[
 R_n=-p_1R_{n-1}-p_2R_{n-2}-\cdots-p_nR_0.
\]
По линейности все слагаемые снова имеют повторённую строку, кроме $-p_1R_{n-1}$. Следовательно, $W'=-p_1W$.
\item Чтобы не делить на потенциально нулевой $W$, положим $P(t)=\int_{t_0}^t p_1(s)\,ds$. Тогда
\[
 (e^{P(t)}W(t))'=e^{P(t)}(p_1W+W')=0.
\]
Значит, $e^{P(t)}W(t)=W(t_0)$, что и даёт формулу.
\end{enumerate}\finish
\remember{Формула верна и при $W(t_0)=0$. Вронскиан либо всюду нулевой, либо нигде не нулевой на интервале.}

\topic{30}{Необходимое условие Эйлера}{§ 3.4, с. 76--78}
\program{Простейшая задача вариационного исчисления. Необходимые условия локального экстремума.}
\claim Пусть $F\in C^2$ в окрестности графика $(t,y_*(t),y_*'(t))$, $y_*\in C^2[a,b]$ --- слабый локальный экстремум
\[
 J[y]=\int_a^bF(t,y(t),y'(t))\,dt
\]
с фиксированными $y(a),y(b)$. Тогда на $(a,b)$
\[
 F_y(t,y_*,y_*')-\frac d{dt}F_p(t,y_*,y_*')=0,\qquad p=y'.
\]
\idea{Изменять функцию вдоль произвольной вариации; одномерная производная в точке экстремума равна нулю.}
\uses{Ферма; дифференцирование интеграла; интегрирование по частям; основную лемму вариационного исчисления}
\begin{enumerate}
\item Берём произвольную $\eta\in C^1[a,b]$ с $\eta(a)=\eta(b)=0$. Функции $y_s=y_*+s\eta$ имеют те же концы и при малом $|s|$ близки к $y_*$ в норме $C^1$.
\item Поэтому $j(s)=J[y_s]$ имеет экстремум при $s=0$, и $j'(0)=0$. Гладкость $F$ и компактность отрезка позволяют дифференцировать под интегралом:
\[
 0=\int_a^b(F_y\eta+F_p\eta')\,dt.
\]
Все производные $F$ вычисляются вдоль $y_*$.
\item Интегрируем второй член по частям. Граничный член $[F_p\eta]_a^b$ равен нулю, поэтому
\[
 \int_a^b\left(F_y-\frac d{dt}F_p\right)\eta\,dt=0
\]
для каждой допустимой $\eta$.
\item Основная лемма говорит: если непрерывная $g$ удовлетворяет $\int g\eta=0$ для всех таких вариаций, то $g=0$. Иначе в малой окрестности точки с $g>0$ (или $g<0$) можно взять неотрицательную ненулевую гладкую вариацию с носителем внутри неё; интеграл имел бы строгий знак. Применяем это к скобке.
\end{enumerate}\finish
\remember{Уравнение Эйлера необходимо, но само по себе не гарантирует экстремума.}

\topic{31}{Формула Байеса}{§ 4.1.2--4.1.3, с. 79--80}
\program{Полная система событий. Формула полной вероятности. Формула Байеса. Независимость событий и классов событий.}
\claim Пусть $(B_i)$ --- конечное или счётное разбиение пространства исходов на события с $P(B_i)>0$, а $P(A)>0$. Тогда
\[
 P(B_k\mid A)=\frac{P(A\mid B_k)P(B_k)}{\sum_iP(A\mid B_i)P(B_i)}.
\]
\idea{Один и тот же совместный исход $A\cap B_k$ записать через две условные вероятности, а знаменатель разложить по гипотезам.}
\uses{определение условной вероятности; счётную аддитивность}
\begin{enumerate}
\item По определению
\[
 P(B_k\mid A)=\frac{P(A\cap B_k)}{P(A)}.
\]
Условие $P(A)>0$ нужно для деления.
\item Поскольку $P(B_k)>0$, другое условное выражение даёт
\[
 P(A\cap B_k)=P(A\mid B_k)P(B_k).
\]
\item События $A\cap B_i$ попарно не пересекаются, а их объединение равно $A$, потому что $B_i$ образуют разбиение. По аддитивности
\[
 P(A)=\sum_iP(A\cap B_i)=\sum_iP(A\mid B_i)P(B_i).
\]
Это одновременно вывод формулы полной вероятности.
\item Подставляя числитель из шага 2 и знаменатель из шага 3 в шаг 1, получаем формулу Байеса.
\end{enumerate}\finish
\remember{Независимость событий здесь не предполагается. Формула пересчитывает вероятности гипотез после наблюдения события $A$.}

\topic{32}{Ожидание и дисперсия нормального распределения}{§ 4.2.2, с. 81--82}
\program{Математическое ожидание и дисперсия случайной величины, их свойства. Вычисление для нормального распределения.}
\claim Если $X$ имеет нормальную плотность
\[
 p(x)=\frac1{\sigma\sqrt{2\pi}}e^{-(x-a)^2/(2\sigma^2)},\qquad\sigma>0,
\]
то $\E X=a$ и $\Var X=\sigma^2$.
\idea{Заменой переменной перейти к стандартной нормальной величине; первое ожидание зануляется нечётностью, второе вычисляется по частям.}
\uses{формулы ожидания по плотности; замену переменной; интегрирование по частям; нормировку гауссовой плотности}
\begin{enumerate}
\item Положим $Z=(X-a)/\sigma$. Замена $x=a+\sigma z$ показывает, что плотность $Z$ равна $\varphi(z)=e^{-z^2/2}/\sqrt{2\pi}$.
\item Функция $z\varphi(z)$ нечётна и абсолютно интегрируема (экспонента подавляет степень). Поэтому
\[
 \E Z=\int_{-\infty}^{\infty}z\varphi(z)\,dz=0.
\]
\item Имеем $\varphi'(z)=-z\varphi(z)$. На $[-A,A]$ интегрирование по частям даёт
\[
 \int_{-A}^Az^2\varphi(z)\,dz
 =-\int_{-A}^Az\varphi'(z)\,dz
 =-[z\varphi(z)]_{-A}^A+\int_{-A}^A\varphi(z)\,dz.
\]
При $A\to\infty$ граничный член стремится к нулю, поскольку $A e^{-A^2/2}\to0$, а последний интеграл стремится к 1 по нормировке плотности. Значит, $\E Z^2=1$.
\item Из $X=a+\sigma Z$ получаем $\E X=a+\sigma\E Z=a$. Кроме того,
\[
 \Var X=\E(X-a)^2=\sigma^2\E Z^2=\sigma^2.
\]
\end{enumerate}\finish
\remember{Параметры $a$ и $\sigma^2$ становятся ожиданием и дисперсией именно в результате вычисления; их нельзя просто объявить таковыми вместо доказательства.}

\topic{33}{Закон больших чисел Чебышёва}{§ 4.3, с. 82--83}
\program{Неравенство Чебышева и закон больших чисел.}
\claim Пусть $X_j$ независимы, $\E X_j=a$, $\Var X_j\le C<\infty$ для всех $j$. Тогда среднее $\overline X_n=n^{-1}\sum_{j=1}^nX_j$ сходится к $a$ по вероятности. Точнее,
\[
 P(|\overline X_n-a|\ge\varepsilon)\le\frac{C}{n\varepsilon^2}\longrightarrow0
 \quad(\varepsilon>0).
\]
\idea{Дисперсия среднего убывает как $1/n$; неравенство Чебышёва превращает это в малую вероятность отклонения.}
\uses{линейность ожидания; независимость; монотонность ожидания}
\begin{enumerate}
\item \textbf{Сначала выведем Чебышёва.} Для любой $Y$ с конечной дисперсией и $A=\{|Y-\E Y|\ge\varepsilon\}$ имеем поточечную оценку
\[
 (Y-\E Y)^2\ge\varepsilon^2\mathbf1_A.
\]
Ожидания дают $\Var Y\ge\varepsilon^2P(A)$, то есть $P(A)\le\Var Y/\varepsilon^2$.
\item По линейности $\E\overline X_n=a$. Для центрированных $U_j=X_j-a$ независимость даёт $\E U_iU_j=(\E U_i)(\E U_j)=0$ при $i\ne j$. Поэтому при раскрытии квадрата суммы смешанные члены исчезают:
\[
 \Var\overline X_n
 =\frac1{n^2}\sum_{j=1}^n\Var X_j\le\frac Cn.
\]
\item Применяем Чебышёва из шага 1 к $Y=\overline X_n$. Получаем заявленную оценку. При фиксированном $\varepsilon>0$ её правая часть стремится к нулю, что и означает сходимость по вероятности.
\end{enumerate}\finish
\remember{Доказаны обе главные части пункта: само неравенство Чебышёва и вывод из него закона больших чисел. Одинаковая распределённость здесь не требуется.}

\topic{34}{Центральная предельная теорема}{§ 4.4, с. 83--84}
\program{Центральная предельная теорема для независимых одинаково распределённых случайных величин с конечной дисперсией.}
\claim Для независимых одинаково распределённых $X_j$ с $\E X_j=a$ и $0<\Var X_j=\sigma^2<\infty$
\[
 Z_n=\frac{\sum_{j=1}^nX_j-na}{\sigma\sqrt n}
 \xrightarrow{\ d\ }N(0,1).
\]
\idea{Характеристическая функция одного малого слагаемого имеет вид $1-t^2/(2n)+o(1/n)$; независимость превращает сумму в $n$-ю степень.}
\uses{теорему о мажорируемой сходимости; теорему непрерывности Леви; характеристическую функцию $N(0,1)$, равную $e^{-t^2/2}$}
\begin{enumerate}
\item Стандартизуем: $Y_j=(X_j-a)/\sigma$. Тогда $\E Y_j=0$, $\E Y_j^2=1$. Обозначим $\phi(u)=\E e^{iuY_1}$.
\item Для вещественного $v$ положим $r(v)=e^{iv}-1-iv+v^2/2$. По Тейлору $r(v)/v^2\to0$ при $v\to0$; кроме того, $|r(v)|\le v^2$ (остаток после линейного члена не больше $v^2/2$).
\item Следовательно, при $u\to0$ величина $r(uY_1)/u^2\to0$ для каждого исхода и по модулю не больше $Y_1^2$. По мажорируемой сходимости её ожидание стремится к нулю. Итак,
\[
 \phi(u)=1-\frac{u^2}2+o(u^2).
\]
\item По независимости для фиксированного $t$
\[
 \phi_{Z_n}(t)=\bigl[\phi(t/\sqrt n)\bigr]^n
 =\left[1-\frac{t^2}{2n}+o(1/n)\right]^n\longrightarrow e^{-t^2/2}.
\]
Последний предел следует из $\log(1+z)=z+o(z)$ около нуля: логарифм степени стремится к $-t^2/2$.
\item Предельная функция --- характеристическая функция стандартного нормального закона, непрерывная в нуле. По теореме Леви получаем сходимость по распределению.
\end{enumerate}\finish
\remember{Не нужен бесконечный ряд по моментам: второго момента достаточно для одного разложения с остатком $o(u^2)$. Дисперсия должна быть строго положительной.}

\topic{35}{Критерий Коши--Римана}{§ 5.1.1--5.1.2, с. 85--86}
\program{Дифференцируемость функций комплексного переменного. Условия Коши -Римана. Интегральная теорема Коши.}
\claim Для $f=u+iv$ комплексная дифференцируемость в $z_0=x_0+iy_0$ равносильна вещественной дифференцируемости $u,v$ в $(x_0,y_0)$ и равенствам
\[
 u_x=v_y,\qquad u_y=-v_x.
\]
При этом $f'(z_0)=u_x+iv_x$; все частные производные берутся в данной точке.
\idea{Комплексное умножение --- специальное вещественное линейное отображение; Коши--Риман описывает его матрицу.}
\uses{определения комплексной и вещественной дифференцируемости}
\begin{enumerate}
\item \textbf{Необходимость.} Если $f'(z_0)=\alpha+i\beta$, то для $h=s+it$
\[
 f(z_0+h)-f(z_0)=(\alpha+i\beta)(s+it)+o(|h|).
\]
Вещественная часть линейного члена равна $\alpha s-\beta t$, мнимая --- $\beta s+\alpha t$. Значит, $u,v$ вещественно дифференцируемы и
\[
 u_x=\alpha,\quad u_y=-\beta,\quad v_x=\beta,\quad v_y=\alpha.
\]
Отсюда требуемые равенства.
\item \textbf{Достаточность.} Из вещественной дифференцируемости
\[
 \Delta f=(u_x+iv_x)s+(u_y+iv_y)t+o(|h|).
\]
По Коши--Риману $u_y+iv_y=-v_x+iu_x=i(u_x+iv_x)$. Поэтому
\[
 \Delta f=(u_x+iv_x)(s+it)+o(|h|).
\]
\item При $h\ne0$ делим на $h$: остаток, делённый на $h$, стремится к нулю по модулю. Значит, комплексная производная существует и равна $u_x+iv_x$.
\end{enumerate}\finish
\remember{Равенств между частными производными без вещественной дифференцируемости недостаточно. Её обеспечивает, например, непрерывность частных производных около точки.}

\topic{36}{Интегральная формула Коши}{§ 5.2.1, с. 87--89}
\program{Интегральная формула Коши. Разложение функции, регулярной в окрестности точки, в ряд Тейлора.}
\claim Пусть $D$ --- ограниченная область с простой замкнутой кусочно-гладкой границей, $f$ регулярна в окрестности $\overline D$. Для $z\in D$ при положительном обходе границы
\[
 f(z)=\frac1{2\pi i}\oint_{\partial D}\frac{f(w)}{w-z}\,dw.
\]
\idea{Стянуть контур к маленькой окружности вокруг $z$: на ней $f(w)$ почти постоянно и равно $f(z)$ с малой ошибкой.}
\uses{интегральную теорему Коши для полной границы области с отверстием; непрерывность $f$}
\begin{enumerate}
\item Выберем $r>0$ так, чтобы $\overline B(z,r)\subset D$. Функция $f(w)/(w-z)$ регулярна между исходной границей и окружностью $\gamma_r:|w-z|=r$.
\item По теореме Коши и ориентации отверстия
\[
 \oint_{\partial D}\frac{f(w)}{w-z}\,dw
 =\oint_{\gamma_r}\frac{f(w)}{w-z}\,dw,
\]
где обе записанные кривые обходятся против часовой стрелки.
\item Разложим числитель справа на $f(z)+[f(w)-f(z)]$. По параметризации $w=z+re^{it}$ имеем
\[
 \oint_{\gamma_r}\frac{dw}{w-z}=\int_0^{2\pi}i\,dt=2\pi i.
\]
Первое слагаемое равно $2\pi i f(z)$.
\item Для второго слагаемого оценка интеграла даёт
\[
 \left|\oint_{\gamma_r}\frac{f(w)-f(z)}{w-z}\,dw\right|
 \le2\pi\max_{|w-z|=r}|f(w)-f(z)|\longrightarrow0
\]
при $r\to0$ по непрерывности в $z$.
\item Левый интеграл шага 2 не зависит от $r$. Устремляя $r$ к нулю, получаем его равенство $2\pi if(z)$; деление даёт формулу.
\end{enumerate}\finish
\remember{Интеграл от $1/(w-z)$ равен $2\pi i$, а не нулю: внутри маленькой окружности находится особенность ядра.}

\topic{37}{Теорема Лорана: существование и единственность}{§ 5.3.1, с. 91--92}
\program{Разложение функции, регулярной в кольце, в ряд Лорана. Изолированные особые точки однозначного характера.}
\claim Если $f$ регулярна в кольце $r<|z-a|<R$, то единственным образом
\[
 f(z)=\sum_{n=-\infty}^{\infty}c_n(z-a)^n,
 \qquad c_n=\frac1{2\pi i}\oint_{|w-a|=\rho}\frac{f(w)}{(w-a)^{n+1}}\,dw,
\]
где $r<\rho<R$. Ряд абсолютно и равномерно сходится на каждом замкнутом подкольце; окружности ориентированы положительно.
\idea{Применить Коши между двумя окружностями. Внешняя окружность даст неотрицательные степени, внутренняя --- отрицательные.}
\uses{формулу Коши для области с отверстием; геометрический ряд; почленное интегрирование равномерно сходящихся рядов}
\begin{enumerate}
\item Зафиксируем подкольцо $\alpha\le|z-a|\le\beta$, где $r<\alpha<\beta<R$, и выберем $r<\rho_1<\alpha<\beta<\rho_2<R$. Обозначим окружности этих радиусов через $\gamma_1,\gamma_2$.
\item Формула Коши на области между ними даёт
\[
 f(z)=\frac1{2\pi i}\oint_{\gamma_2}\frac{f(w)}{w-z}\,dw
 -\frac1{2\pi i}\oint_{\gamma_1}\frac{f(w)}{w-z}\,dw.
\]
Минус возникает из ориентации внутренней границы.
\item На внешней окружности разложим
\[
 \frac1{w-z}=\sum_{k=0}^\infty\frac{(z-a)^k}{(w-a)^{k+1}},
 \qquad\frac{|z-a|}{|w-a|}\le\frac\beta{\rho_2}<1.
\]
На внутренней окружности разложим с учётом минуса:
\[
 -\frac1{w-z}=\sum_{k=0}^\infty\frac{(w-a)^k}{(z-a)^{k+1}},
 \qquad\frac{|w-a|}{|z-a|}\le\frac{\rho_1}\alpha<1.
\]
Оба ряда равномерны на выбранных контурах и подкольце, поэтому их можно интегрировать почленно. Получаем соответственно степени $n\ge0$ и $n=-k-1<0$ с заявленными коэффициентами.
\item Для каждого целого $n$ функция $f(w)/(w-a)^{n+1}$ регулярна в кольце. По теореме Коши её интегралы по любым двум концентрическим окружностям равны: интеграл по полной границе промежуточного кольца нулевой. Поэтому коэффициенты не зависят от $\rho$.
\item Если $M_j=\max_{\gamma_j}|f|$, то для положительной части $|c_k(z-a)^k|\le M_2(\beta/\rho_2)^k$. Для отрицательной части
\[
 |c_{-k-1}(z-a)^{-k-1}|\le M_1(\rho_1/\alpha)^{k+1}.
\]
Геометрические мажоранты дают абсолютную и равномерную сходимость.
\item \textbf{Единственность.} Умножим любой такой ряд на $(z-a)^{-m-1}$ и почленно проинтегрируем по окружности внутри кольца. Интеграл $\oint(z-a)^jdz$ равен нулю при целом $j\ne-1$ и равен $2\pi i$ при $j=-1$. Останется только $2\pi i c_m$. Значит, любой ряд обязан иметь указанные коэффициенты.
\end{enumerate}\finish
\remember{В отличие от доказательства только формулы коэффициентов, здесь доказано само существование разложения. Два геометрических ряда отвечают двум частям Лорана.}

\topic{38}{Теорема о вычетах}{§ 5.4.1--5.4.2, с. 94--96}
\program{Вычеты. Вычисление интегралов по замкнутому контуру при помощи вычетов.}
\claim Пусть $D$ --- ограниченная область с простой замкнутой кусочно-гладкой границей, а $f$ регулярна в окрестности $\overline D$, кроме конечного числа изолированных особых точек $a_1,\ldots,a_m\in D$. Тогда при положительном обходе
\[
 \oint_{\partial D}f(z)\,dz=2\pi i\sum_{j=1}^m\operatorname{res}_{a_j}f.
\]
\idea{Вырезать маленькие диски вокруг особенностей. В оставшейся области применить Коши; интегралы по маленьким окружностям выделят вычеты.}
\uses{интегральную теорему Коши для полной границы области с отверстиями; ряд Лорана; раздел 37}
\begin{enumerate}
\item Выбираем попарно непересекающиеся замкнутые диски $\overline B(a_j,r_j)\subset D$, каждый из которых содержит только одну особую точку. Это возможно, поскольку точек конечное число и все они внутренние.
\item Обозначим их окружности через $\gamma_j$, ориентированные против часовой стрелки, и вырежем диски из $D$. На окрестности замыкания оставшейся области $f$ регулярна.
\item Её положительная граница состоит из внешней $\partial D$ и окружностей отверстий, обходимых по часовой стрелке. Поэтому теорема Коши даёт
\[
 0=\oint_{\partial D}f(z)\,dz-\sum_{j=1}^m\oint_{\gamma_j}f(z)\,dz.
\]
\item Вокруг $a_j$ есть ряд Лорана. Вычет --- его коэффициент $c_{-1}$. По формуле коэффициентов (раздел 37, $k=-1$)
\[
 \oint_{\gamma_j}f(z)\,dz=2\pi i\,c_{-1}
 =2\pi i\operatorname{res}_{a_j}f.
\]
\item Подставляем в шаг 3 и переносим сумму вправо. Получаем утверждение.
\end{enumerate}\finish
\remember{Минус перед внутренними окружностями возникает из ориентации отверстий. Особенностей на исходном контуре быть не должно.}

\end{document}
